\section{The same original projective twists add under tensor product}

\begin{lemma}[The same original projective twists add]
\label{lem:coherent-projective-original-twist-addition}
\uses{thm:scheme-original-frame-gluing,lem:scheme-original-tensor-frame-transitions}
\lean{EnvelopeProperty.ProjectiveTwistTensorAddition.tensorIso}
\leanok
For the genuine coordinate-hyperplane presentations of the usual twists
on polynomial projective space, $\mathcal O_j(a)\otimes\mathcal O_k(b)$
is isomorphic to $\mathcal O_\ell(a+b)$ for every choice of coordinate
indices $j,k,\ell$ and integer degrees $a,b$.
\end{lemma}
\begin{proof}
On each original standard coordinate open, use the original coefficient-one
frames. On an overlap the degree-$a$ and degree-$b$ frames have transition
factors $\tau^a$ and $\tau^b$, independent of the chosen hyperplane index.
The tensor frame therefore has transition $\tau^{a+b}$, exactly that of
the original target frame. The framewise coefficient-preserving maps
thus agree on overlaps, and the actual module-Hom sheaf glues them to
an original whole morphism. The framewise inverse maps glue in the same
way. Each composite is the identity on every member of the standard
cover, hence is the identity globally by separatedness of the same
module-Hom sheaf. These two original morphisms give the isomorphism.
\end{proof}

\begin{lemma}[Tensor addition preserves the standard frames]
\label{lem:projective-twist-tensor-addition-frame}
\uses{lem:coherent-projective-original-twist-addition,thm:scheme-original-frame-gluing,lem:scheme-original-tensor-frame-transitions}
\lean{EnvelopeProperty.ProjectiveTwistTensorAddition.tensorIso_restrict}
\leanok
The tensor-addition isomorphism may be chosen so that, on each standard
coordinate open, it is exactly the comparison from the tensor of the
degree-$a$ and degree-$b$ coefficient-one frames to the degree-$(a+b)$
coefficient-one frame. These choices are made simultaneously for every
choice of coordinate presentations and integer degrees.
\end{lemma}
\begin{proof}
Use the isomorphism obtained by gluing the specified frame comparisons
in the preceding proof. The restriction identity for gluing identifies
its restriction to each member of the cover with the prescribed local
comparison. Thus the tensor frame and the target frame have identical
coefficients under this restriction, for all the indicated parameters.
\end{proof}
